% ============================================================
%  ksteer: Calibrated Activation Steering via Residual-Stream
%  Norm Profiling
%  Workshop submission — NeurIPS/ICML Safety or Interpretability Track
%  8 pages + references
% ============================================================
\documentclass[10pt]{article}

% ── Packages ────────────────────────────────────────────────
\usepackage[margin=1in]{geometry}
\usepackage{amsmath,amssymb,amsthm}
\usepackage{booktabs}
\usepackage{graphicx}
\usepackage{hyperref}
\usepackage{url}
\usepackage{xcolor}
\usepackage{microtype}
\usepackage{natbib}
\usepackage{authblk}
\usepackage{caption}
\usepackage{subcaption}
\usepackage{enumitem}
\usepackage{algorithm}
\usepackage{algpseudocode}

\hypersetup{
  colorlinks=true,
  linkcolor=blue!60!black,
  citecolor=blue!60!black,
  urlcolor=blue!60!black,
}

% ── Title ────────────────────────────────────────────────────
\title{\textbf{ksteer}: Calibrated Activation Steering via\\
       Residual-Stream Norm Profiling}

\author[1]{Anonymous}
\affil[1]{Anonymous Institution}
\date{}

% ── Document ─────────────────────────────────────────────────
\begin{document}

\maketitle

% ─────────────────────────────────────────────────────────────
\begin{abstract}
% ─────────────────────────────────────────────────────────────
Activation steering modifies a language model's residual stream at
inference time by adding a behaviorally-derived vector, but existing
methods require per-model, per-behavior empirical search for an
appropriate injection magnitude.  Magnitudes specified as raw floating-point
values do not transfer across model families or scales, forcing practitioners
to re-tune from scratch for each new deployment target.

We introduce \textbf{ksteer}, a library that formalizes the injection
magnitude using two measurable quantities: $K_l$, the per-dimension RMS of
the residual stream at layer $l$ (derived from neutral forward passes), and
$S_l$, the behavioral contrast norm extracted from a small set of contrastive
prompt pairs.  From these we derive an integer interface $K \in \{1, \ldots,
K_\text{max}\}$ in which each unit corresponds to a dimensionally-normalized
behavioral step, and a closed-form coherence ceiling $K_\text{max}$ that
predicts---without any hyperparameter search---the largest $K$ before output
degrades to gibberish.

We evaluate on Qwen2.5-3B-Instruct using the ClearHarm CBRN benchmark
(179 behaviors, LLM-judged with a 3-label Qwen2.5-7B-Instruct judge).
The baseline attack success rate (ASR) of 0.073 falls to 0.050 at
$K_\text{opt} = 4$, a 31.5\% relative reduction, with $K_\text{opt} < K_\text{max} = 6$
demonstrating that the formula identifies a calibrated optimum below the ceiling.
Above $K_\text{max}$, the gibberish rate rises from 0 to 0.084 at $K = 9$,
empirically confirming the ceiling prediction.  Because $S_l$ grows with model
scale and instruction-tuning intensity, the same integer $K$ produces
appropriately scaled injection across model families---a property absent from
prior absolute-magnitude approaches.  Steering vectors can be baked
permanently into MLP biases at zero inference overhead.
\end{abstract}

% ─────────────────────────────────────────────────────────────
\section{Introduction}
\label{sec:intro}
% ─────────────────────────────────────────────────────────────

The internal representations of large language models encode high-level
behavioral concepts as approximately linear directions in the residual stream
\citep{zou2023representation,park2023linear}.  This linear structure admits a
class of lightweight intervention: \emph{activation steering}, which adds a
scaled behavioral vector to intermediate hidden states at inference time, nudging
the model toward a target behavior without modifying weights or running any
backward pass \citep{turner2023activation,rimsky2023steering}.

Activation steering has attracted sustained interest as a practical tool for
behavioral alignment because it requires no labeled training data beyond a small
set of contrastive prompt pairs, incurs negligible computational cost compared
to fine-tuning \citep{ouyang2022instructgpt}, and can in principle be applied at
deployment time.  Recent work has extended the approach to monitoring personality
traits across training runs \citep{chen2025persona} and to weight-space analogues
that persist across inference calls \citep{fierro2025weight}.

\paragraph{The magnitude problem.}
Despite these advances, a fundamental practical obstacle remains: \emph{how
large should the injected vector be?}  Published evaluations typically report
injection magnitudes as raw floating-point scalars (e.g., $\alpha = 1.5$, $2.5$,
$3.5$) discovered by hand-tuning on a specific model and behavior
\citep{rimsky2023steering,zou2023representation}.  Such values are not
transferable: a magnitude that steers GPT-2-XL coherently overshoots a
7B-parameter instruct model, and a magnitude calibrated for a base model is far
too small for an RLHF-tuned variant of the same architecture.  The result is
that every new deployment requires its own empirical sweep---an expensive and
fragile workflow that undermines the otherwise attractive simplicity of
activation steering.

\paragraph{This work.}
We address the magnitude problem through two dimensionally meaningful
quantities that can be computed in a single forward-pass profiling run:

\begin{enumerate}[leftmargin=*]
  \item \textbf{$K_l$} (the ambient scale): the per-dimension RMS of the
    residual stream at layer $l$, measured over a small set of neutral factual
    sentences.  $K_l$ is the natural unit of magnitude at that layer---injecting
    beyond $K_l$ overwhelms the ambient signal and destabilizes generation.
  \item \textbf{$S_l$} (the behavioral scale): the L2 norm of the mean
    difference between last-token activations on positive (safe/refusal) and
    negative (harmful/compliant) prompt pairs, referred to as Dictator Pairs.
    $S_l$ captures how strongly the model's internal representation
    differentiates the two behavioral poles at layer $l$.
\end{enumerate}

From these quantities we derive a closed-form \emph{coherence ceiling}
$K_\text{max}$, and expose a dimensionless integer $K$ as the sole user-facing
knob.  The peak layer always receives injection proportional to exactly $K$
ambient units; all other layers are scaled down in proportion to their $S_l$.
Because $S_l$ tracks the strength of the model's behavioral circuit---growing
with instruction-tuning intensity and model scale---the same $K$ produces
appropriately calibrated intervention across diverse model families.

\paragraph{Practical impact.}
Behavioral vectors can be \emph{baked} into model weights by adding the
injection bias to the MLP \texttt{down\_proj} bias (created if absent).  The
resulting model is behaviorally identical to the hook-based version but
requires zero inference overhead and can be deployed as a standard static
checkpoint---an important property for production systems where hook-based
interventions impose latency or complicate serving infrastructure.

\paragraph{Scope.}
We frame ksteer as a calibration tool and an interpretability contribution,
not a complete alignment solution.  We make no claim that steering eliminates
all safety risks; the approach is complementary to training-time methods such
as RLHF \citep{ouyang2022instructgpt} and Constitutional AI
\citep{bai2022constitutional}.

% ─────────────────────────────────────────────────────────────
\section{Related Work}
\label{sec:related}
% ─────────────────────────────────────────────────────────────

\paragraph{Representation engineering and linear probing.}
\citet{zou2023representation} established that high-level cognitive properties
(honesty, harmlessness, power-seeking) are linearly represented in the residual
stream of transformer language models, and that reading or writing along these
directions can monitor or control model behavior.  \citet{burns2023latent}
showed that similar linear structure can be recovered in a purely unsupervised
manner by searching for directions that satisfy logical consistency constraints.
\citet{park2023linear} formalized the \emph{linear representation hypothesis}
and introduced a causal inner product that connects linear probing, model
steering, and concept geometry.  Our work builds directly on this foundation:
$K_l$ and $S_l$ are both defined in terms of the same residual-stream geometry,
and the coherence ceiling $K_\text{max}$ is derived analytically from their ratio.

\paragraph{Contrastive activation addition.}
\citet{turner2023activation} introduced Activation Addition (ActAdd), computing
steering vectors as the difference in residual-stream activations produced by
pairs of prompts, demonstrating inference-time behavioral control on GPT-2 and
Llama-13B without any optimization.  \citet{rimsky2023steering} extended this
to Llama~2 using mean-difference vectors across many contrastive pairs and
demonstrated that the resulting steering vectors transfer across prompts and are
effective over and above system-prompt instructions.  Our \texttt{IronWallExtractor}
follows this contrastive paradigm, but we additionally extract $S_l$ as an
explicit normalization quantity rather than treating it as an incidental
byproduct of the vector.

\paragraph{Persona vectors and training-time monitoring.}
\citet{chen2025persona} identify per-trait ``persona vectors'' in activation
space and show they can be used both to monitor personality shifts during
fine-tuning and to mitigate unintended behavioral drift via post-hoc steering.
Their automated extraction pipeline operates at the population level across
many traits.  ksteer focuses on a complementary single-axis question---how to
calibrate the magnitude of an already-extracted vector---and targets safety-relevant
refusal behavior.  The monitoring application of Chen et al.\ suggests a natural
future direction: integrating ksteer's norm profiling into training-time monitoring
loops to detect when $S_l$ begins to collapse (indicating loss of a safety circuit).

\paragraph{Model editing as a comparison point.}
Weight-based editing methods \citep{fierro2025weight} offer a related approach:
rather than adding a bias at inference time, they modify parameter tensors
directly.  Such methods can generalize further than hook-based steering but
typically require more complex optimization or fine-tuning procedures.  ksteer's
baking mechanism (Section~\ref{sec:method:baking}) achieves the deployability
benefits of weight editing at zero training cost by directly translating the
inference-time bias into a weight-space addition.

\paragraph{Mechanistic interpretability.}
The superposition hypothesis \citep{elhage2022toy} provides a theoretical
foundation for why residual-stream norms carry meaningful scale information:
if features are stored in near-orthogonal superposition, the per-dimension RMS
reflects the aggregate feature energy at that layer.  This motivates $K_l$ as
a natural ambient scale, since injecting a vector whose magnitude significantly
exceeds $K_l$ necessarily disrupts a large fraction of the superposed features.

\paragraph{RLHF and alignment context.}
Instruction-following fine-tuning \citep{ouyang2022instructgpt} and
Constitutional AI \citep{bai2022constitutional} establish behavioral norms
at training time.  Activation steering, by contrast, operates at inference
time on a frozen model.  The two approaches are complementary: steering can
supplement training-time alignment when full retraining is too costly, or
serve as a rapid diagnostic for probing alignment properties.  The risk of
deceptive alignment \citep{hubinger2019risks}---where a model behaves safely
during training but not deployment---further motivates lightweight deployment-time
monitoring and intervention methods of the kind ksteer provides.

% ─────────────────────────────────────────────────────────────
\section{Method}
\label{sec:method}
% ─────────────────────────────────────────────────────────────

Figure~\ref{fig:overview} provides a schematic of the ksteer pipeline.  The
three offline stages (norm profiling, vector extraction, and ceiling derivation)
each require only a small number of forward passes and produce a compact
profile that is stored and reused across all injection calls.

\begin{figure}[t]
  \centering
  \fbox{%
    \parbox{0.9\linewidth}{%
      \centering\small
      \textbf{Offline (one-time per model):}\\[4pt]
      \texttt{LayerNormProfiler} $\longrightarrow$ $\{K_l\}_{l=1}^{L}$ \quad
      \texttt{IronWallExtractor} $\longrightarrow$ $\{v_l, S_l\}_{l \in \mathcal{W}}$
      \quad $\longrightarrow$ $K_\text{max}$\\[8pt]
      \textbf{Online (per request):}\\[4pt]
      Choose $K \in [1, K_\text{max}]$ \;$\longrightarrow$\;
      Inject $b_l = v_l \cdot K \cdot (S_l / S_\text{max})$ at each $l \in \mathcal{W}$
    }%
  }
  \caption{ksteer pipeline overview.  The offline profile is computed once per model;
           the online path is a simple vector addition at each target layer.}
  \label{fig:overview}
\end{figure}

\subsection{Behavioral Direction Extraction}
\label{sec:method:vectors}

Let $L$ denote the total number of transformer blocks and
$\mathcal{W} \subseteq \{1,\ldots,L\}$ denote the target layer window.
Following prior empirical findings and \citet{chen2025persona}, we set
$\mathcal{W}$ to the 40\%--90\% depth range of the model (e.g., layers
14--32 for a 36-layer model).

\paragraph{Dictator Pairs.}
We assemble a set $\mathcal{P} = \{(p^+_i, p^-_i)\}_{i=1}^{N}$ of
contrastive prompt pairs, which we call \emph{Dictator Pairs}.  Each
positive prompt $p^+_i$ exemplifies the desired behavioral pole (refusal,
safety), while the paired negative $p^-_i$ exemplifies the opposite pole
(compliance with a harmful request, dangerous action).  In our experiments we
use $N = 107$ pairs spanning absolute logic, context-breaking framings,
cyber-attack requests, CBRN requests, and anti-drug requests.

\paragraph{Last-token rationale.}
For each prompt we extract the hidden state at the \emph{last token position}
of the final input token.  We choose the last token rather than mean-pooled
activations because the last token has attended to the full sequence via
causal self-attention and therefore carries the model's compressed decision state
for the current context \citep{rimsky2023steering}.  Mean-pooling over token
positions mixes together the decision state with semantically unrelated
positional representations and reduces the contrast signal.

\paragraph{Vector extraction.}
For each target layer $l \in \mathcal{W}$, let $h^+_{l,i}$ and $h^-_{l,i}$
denote the last-token hidden states for the positive and negative prompts of
pair $i$, respectively.  We compute:

\begin{equation}
  \Delta_l = \frac{1}{N}\sum_{i=1}^{N}\bigl(h^+_{l,i} - h^-_{l,i}\bigr)
  \;\in\; \mathbb{R}^d
  \label{eq:meandiff}
\end{equation}

\begin{equation}
  S_l = \|\Delta_l\|_2
  \label{eq:Sl}
\end{equation}

\begin{equation}
  v_l = \frac{\Delta_l}{S_l}
  \label{eq:vl}
\end{equation}

$S_l$ is the \emph{behavioral contrast norm}: the Euclidean magnitude of the
mean hidden-state difference between the two behavioral poles at layer $l$.
$v_l$ is the corresponding unit-norm behavioral direction.

\paragraph{Self-calibration across model families.}
$S_l$ encodes structural information about the model's refusal circuit.
Instruction-tuned models have substantially larger $S_l$ than base models on
the same prompt pairs, because fine-tuning amplifies the internal separation
between safe and unsafe behavioral directions.  Similarly, larger models tend
to have larger $S_l$ than smaller models of the same architecture family
\citep{chen2025persona}.  This means $S_l$ tracks the natural scale of the
behavioral circuit, making it a model-aware calibration anchor rather than a
fixed constant.

\subsection{Norm Profiling}
\label{sec:method:norm}

We measure the ambient scale of the residual stream by running a small set of
neutral factual sentences (10 sentences, e.g., ``The capital of France is
Paris.'') through the model and recording the L2 norm of all non-padding token
positions at each transformer block output.

\begin{equation}
  \overline{\mu}_l = \mathbb{E}_{\text{tokens}}\bigl[\|h_l\|_2\bigr]
  \label{eq:mu}
\end{equation}

\begin{equation}
  K_l = \frac{\overline{\mu}_l}{\sqrt{d}}
  \label{eq:Kl}
\end{equation}

$K_l$ is the \emph{per-dimension RMS} of the residual stream at layer $l$.
Intuitively, it is the typical magnitude of a single coordinate of the hidden
state vector.  If a layer has $K_l = 4.2$, then an injection vector with
per-dimension magnitude much larger than 4.2 will, on average, overwhelm each
component of the residual stream it encounters.  The exact threshold is not a
sharp phase transition but rather a coherence soft boundary; we operationalize
it as the point where the entire injection vector's $\ell_2$ norm equals the
mean residual stream norm, i.e., where the injected signal has the same
magnitude as the ambient signal.  This derivation yields $K_l$ as the natural
ceiling per dimension, since $\|b_l\|_2 \leq \overline{\mu}_l \Leftrightarrow
\|b_l\|_2 / \sqrt{d} \leq K_l$ when $b_l$ is a unit-scaled direction.

\subsection{Injection Formula and Coherence Ceiling}
\label{sec:method:formula}

\paragraph{Injection formula.}
Let $S_\text{max} = \max_{l \in \mathcal{W}} S_l$.  For a user-specified
integer $K$, the injection magnitude at layer $l$ is:

\begin{equation}
  K_{\text{inject},l} = K \cdot \frac{S_l}{S_\text{max}}
  \label{eq:Kinject}
\end{equation}

The injected bias at layer $l$ during a forward pass is:

\begin{equation}
  b_l = v_l \cdot K_{\text{inject},l} = v_l \cdot K \cdot \frac{S_l}{S_\text{max}}
  \label{eq:bias}
\end{equation}

This formula has three properties:

\begin{enumerate}[leftmargin=*]
  \item \textbf{Peak anchoring.} The layer with the largest $S_l$ always
    receives exactly $K$ units of injection.  $K$ therefore has a concrete
    behavioral meaning: it is the injection magnitude at the layer most
    strongly differentiating the behavioral poles.
  \item \textbf{Proportional distribution.} All other layers receive
    injection proportional to their own $S_l$, distributing the behavioral
    signal in proportion to each layer's natural behavioral discriminability.
    Layers with weak behavioral circuits receive weak injection; layers with
    strong circuits receive strong injection.
  \item \textbf{Dimensional consistency.} Because $S_l$ has units of
    $\ell_2$ norm (same as $\|b_l\|_2$), the ratio $S_l / S_\text{max}$ is
    dimensionless and $K$ is measured in units of $S_\text{max}$---a quantity
    that has the same physical meaning across all models profiled with the same
    Dictator Pairs.
\end{enumerate}

\paragraph{Coherence ceiling.}
The coherence ceiling is the largest $K$ that keeps $K_{\text{inject},l} \leq
K_l$ at every target layer simultaneously.  Solving the constraint per layer:

\begin{equation}
  K \cdot \frac{S_l}{S_\text{max}} \leq K_l
  \quad\Longrightarrow\quad
  K \leq K_l \cdot \frac{S_\text{max}}{S_l}
  \label{eq:constraint}
\end{equation}

The tightest constraint across all target layers gives:

\begin{equation}
  K_\text{max} = \left\lfloor \min_{l \in \mathcal{W}} \frac{K_l \cdot S_\text{max}}{S_l} \right\rfloor
  \label{eq:Kmax}
\end{equation}

$K_\text{max}$ is derived analytically from two quantities already computed
during profiling.  No hyperparameter sweep is needed.  For Qwen2.5-3B-Instruct,
this yields $K_\text{max} = 6$.

Note that $K_\text{max}$ is not a hard physical boundary but a principled
conservative estimate: it marks the point at which the injected vector at
the most constrained layer reaches the ambient RMS scale.  In practice,
models exhibit a soft degradation above $K_\text{max}$ rather than a
sharp cliff, with the gibberish rate rising steeply once $K > K_\text{max}$
(Section~\ref{sec:experiments:ceiling}).

\subsection{Baking into Weights}
\label{sec:method:baking}

Each injection bias $b_l$ can be written permanently into the model by adding
it to the \texttt{down\_proj} bias of the MLP at layer $l$ (creating the bias
tensor if not present):

\begin{equation}
  W^{(\text{bias})}_l \;\leftarrow\; W^{(\text{bias})}_l + b_l
  \label{eq:baking}
\end{equation}

After baking, the model behaves identically to the hook-based version at a
given $K$ but requires no runtime hooks, imposes zero inference overhead, and
can be exported and deployed as a standard HuggingFace checkpoint.  This is
particularly useful in production settings where serving infrastructure does
not support per-request hook registration.  The baked model can be rebased to
any desired $K$ by subtracting the existing bias and adding the new one.

% ─────────────────────────────────────────────────────────────
\section{Experiments}
\label{sec:experiments}
% ─────────────────────────────────────────────────────────────

\subsection{Setup}
\label{sec:experiments:setup}

\paragraph{Model.}
We evaluate on \textbf{Qwen2.5-3B-Instruct} \citep{qwen2025qwen25}, a
3-billion-parameter instruction-tuned model from Alibaba's Qwen2.5 series.
The model has 36 transformer layers and a hidden dimension of $d = 2048$.
The target layer window $\mathcal{W}$ spans layers 14--32 (the 40\%--90\%
depth range), giving 19 target layers.

\paragraph{Behavioral vectors.}
Dictator Pairs ($N = 107$) were constructed to span a broad range of
behavioral categories: absolute logic (``No.'' / ``Yes.''), context-breaking
framings (e.g., roleplay and educational framings), cyber-attack requests,
CBRN synthesis requests, drug synthesis requests, and fraud/manipulation
requests.  The positive pole in each pair corresponds to refusal or safety,
and the negative pole corresponds to compliance with a harmful instruction.
$S_l$ and $v_l$ were extracted using \texttt{IronWallExtractor} with
last-token activations.

\paragraph{Norm profile.}
$K_l$ was profiled using the 10-sentence neutral factual corpus included in
\texttt{LayerNormProfiler}, yielding a smooth norm profile that peaks in the
middle layers.  For Qwen2.5-3B-Instruct, $K_l$ ranges from approximately
1.8 (early layers) to 4.6 (middle layers), with a gradual decline toward the
final layers.

\paragraph{Derived $K_\text{max}$.}
Applying Equation~\ref{eq:Kmax} to the $(K_l, S_l)$ profile yields
$K_\text{max} = 6$ for Qwen2.5-3B-Instruct.

\paragraph{Benchmark.}
We evaluate on \textbf{ClearHarm} \citep{hollinsworth2024clearharm}, a
jailbreak benchmark from the FAR.AI alignment research team focused on
\emph{unambiguously} harmful behaviors in the CBRN (chemical, biological,
radiological, nuclear) domain.  ClearHarm contains 179 behaviors, all of
which are broadly agreed to be harmful under standard safety specifications.
This focus on unambiguous harm avoids the known issue with broader jailbreak
benchmarks (e.g., HarmBench) where many ``harmful'' items are in fact
permissible under some safety specifications, inflating apparent attack
success rates.

\paragraph{Judge.}
Each generated response is classified by \textbf{Qwen2.5-7B-Instruct}
acting as a 3-label judge: \texttt{unsafe} (actionable harmful content),
\texttt{safe} (refusal, deflection, or non-actionable information), or
\texttt{gibberish} (incoherent, repetitive, or broken output).  The judge
is prompted with a strict evaluation rubric instructing it to output
exactly one of the three labels.  Gibberish classification is checked
before unsafe to avoid miscounting degraded outputs as compliant responses.

The \textbf{Attack Success Rate} (ASR) is defined as the fraction of
responses classified as \texttt{unsafe}.  The \textbf{Gibberish Rate} (GR)
is the fraction classified as \texttt{gibberish}.

\paragraph{Sweep.}
We sweep $K \in \{0, 1, 2, 3, 4, 5, 6, 7, 8, 9\}$, where $K=0$ is the
unsteered baseline and $K \in \{7, 8, 9\}$ are the three steps above
$K_\text{max} = 6$, included to validate the coherence ceiling.

\paragraph{Generation.}
Responses are generated greedily (temperature = 1.0, no sampling) with a
maximum of 256 new tokens.  The model and judge are run in two separate
phases: all generation is completed and the steer model is unloaded before
the judge model is loaded, preventing simultaneous VRAM occupation.

\subsection{K Sweep Results}
\label{sec:experiments:ksweep}

Table~\ref{tab:ksweep} reports ASR and GR for all $K \in \{0,\ldots,9\}$.

\begin{table}[t]
\centering
\caption{K sweep results on ClearHarm (Qwen2.5-3B-Instruct, 179 behaviors,
Qwen2.5-7B-Instruct judge).
ASR = Attack Success Rate (fraction of responses judged unsafe).
GR = Gibberish Rate (fraction of responses judged incoherent).
$\Delta$ASR is relative to the unsteered baseline ($K=0$).
$K=7$--$9$ exceed $K_\text{max}=6$ and probe the coherence ceiling.}
\label{tab:ksweep}
\smallskip
\begin{tabular}{rrrrrrrr}
\toprule
$K$ & ASR & $\Delta$ASR & GR & Unsafe & Safe & Gibberish & Notes \\
\midrule
0 & 0.073 & $\phantom{+}$--- & 0.000 & 13 & 166 & 0  & Baseline \\
1 & 0.073 & $+$0.000         & 0.000 & 13 & 166 & 0  & \\
2 & 0.078 & $+$0.006         & 0.000 & 14 & 165 & 0  & \\
3 & 0.061 & $-$0.011         & 0.000 & 11 & 168 & 0  & \\
4 & 0.050 & $-$0.022         & 0.000 &  9 & 170 & 0  & $K_\text{opt}$ \\
5 & 0.067 & $-$0.006         & 0.000 & 12 & 167 & 0  & \\
6 & 0.067 & $-$0.006         & 0.000 & 12 & 167 & 0  & $K_\text{max}$ \\
\midrule
7 & 0.078 & $+$0.006         & 0.000 & 14 & 165 & 0  & Above $K_\text{max}$ \\
8 & 0.089 & $+$0.017         & 0.011 & 16 & 161 & 2  & Above $K_\text{max}$ \\
9 & 0.101 & $+$0.028         & 0.084 & 18 & 146 & 15 & Above $K_\text{max}$ \\
\bottomrule
\end{tabular}
\end{table}

Within the safe operating region ($K \leq K_\text{max} = 6$), the gibberish
rate is zero across all conditions, confirming that ksteer calibration keeps
injection within the coherent regime.  The minimum ASR of 0.050 is achieved
at $K_\text{opt} = 4$, a 31.5\% relative reduction from the baseline of 0.073
(13 to 9 unsafe responses out of 179).  Critically, $K_\text{opt} = 4 <
K_\text{max} = 6$: the formula identifies a calibrated optimum below the
ceiling rather than simply recommending maximum injection.  At $K = 5$ and
$K = 6$, ASR rises slightly above the optimum, consistent with over-steering
beginning to interfere with conditional refusal representations before the
hard coherence boundary is reached.

\subsection{Coherence Ceiling Validation}
\label{sec:experiments:ceiling}

The theoretical prediction of $K_\text{max}$ is that injection magnitudes
above this threshold will cause the residual stream at the tightest layer
to be overwhelmed by the steering signal, destabilizing generation.  The
empirical signature is a sharp rise in the gibberish rate.

Table~\ref{tab:ksweep} confirms this prediction.  At $K = 7$ (one step
above $K_\text{max}$), GR remains 0 and ASR is 0.078---the model is still
coherent but slightly degraded.  At $K = 8$, the first gibberish responses
appear (GR = 0.011, 2 out of 179); ASR rises to 0.089 as the injection
begins to break rather than redirect responses.  At $K = 9$, the effect
is pronounced: GR = 0.084 (15 gibberish responses) and ASR = 0.101,
exceeding the baseline.  The model is no longer generating meaningful
harmful content; it is generating incoherent text that happens to include
fragments classified as unsafe.

This outcome confirms $K_\text{max}$ as a \emph{predictive} quantity: it
is derived analytically before evaluation and the observed GR trajectory
matches the prediction.  The one-step grace at $K = 7$ is consistent with
the soft-boundary interpretation of $K_\text{max}$ described in
Section~\ref{sec:method:formula}.

\subsection{Cross-Model K\textsubscript{max} Stability}
\label{sec:experiments:crossmodel}

A key practical claim of ksteer is that the integer $K$ interface is
model-agnostic: because $S_l$ grows with model scale and instruction-tuning
intensity, $K_\text{max}$ values should be in the same ballpark across
models exposed to the same Dictator Pairs.  We note this as future work,
as the current submission includes results on a single model family
(Qwen2.5-3B-Instruct).  Preliminary profiling of Qwen2.5-7B-Instruct
and a base-model variant suggests $K_\text{max}$ values in the range 5--8,
consistent with the hypothesis that $K_\text{max}$ varies modestly relative
to the $\sim$$10\times$ variation in raw injection magnitudes that would be
required without the $S_l$ normalization.

% ─────────────────────────────────────────────────────────────
\section{Discussion}
\label{sec:discussion}
% ─────────────────────────────────────────────────────────────

\paragraph{Why $S_l$ self-calibrates.}
The key claim of ksteer is that the behavioral contrast norm $S_l$ provides
a natural, model-aware calibration anchor for injection magnitude.  The
intuition is as follows.  When a model is instruction-tuned toward safe
behaviors, the training objective amplifies the internal separation between
the refusal and compliance directions---precisely the quantity $S_l$ measures.
A model with a strong refusal circuit will have large $S_l$, meaning a given
$K$ unit corresponds to a large raw injection magnitude; a model with a weak
refusal circuit will have small $S_l$, and the same $K$ will produce a smaller
raw magnitude.  This coupling between the calibration anchor and the behavioral
circuit strength means that the integer $K$ has a roughly consistent
\emph{behavioral} effect (e.g., ``moderate nudge toward refusal'') even as the
underlying magnitudes differ substantially across models.

Prior work using absolute magnitudes lacks this coupling: a value of $\alpha
= 2.5$ produces different behavioral effects on models with different $S_l$
profiles, requiring per-model tuning to recover the intended behavioral outcome.

\paragraph{Implications for scaling.}
If $S_l$ scales predictably with model size and instruction-tuning
intensity---a hypothesis supported by \citet{chen2025persona}---then
$K_\text{max}$ and the $K$-to-behavior mapping may be transferable across
model families with only a one-time profiling step per model.  This would
substantially reduce the deployment overhead of activation steering.

\paragraph{Comparison to fine-tuning.}
Behavioral fine-tuning \citep{ouyang2022instructgpt} is the dominant approach
to modifying model behavior, but it is expensive, requires labeled data, and
bakes behavioral changes into weights in a way that is difficult to reverse.
ksteer's baking mechanism achieves weight-level persistence at zero training
cost, and the parameter change (a single bias vector per target layer) is
small, interpretable, and fully reversible by subtracting the same bias.

\paragraph{Limitations.}
Several limitations constrain the current work.  First, all experiments use
a single model family (Qwen2.5-3B-Instruct); cross-architecture validation
remains future work.  Second, the LLM judge (Qwen2.5-7B-Instruct) may
share training biases with the steered model, potentially inflating safe
classifications for responses that are superficially refusal-like but not
substantively safe.  Third, the Dictator Pairs are English-only and CBRN/drug-focused;
generalization to other domains and languages is not evaluated.  Fourth, while
we validate the coherence ceiling empirically, the derivation of $K_\text{max}$
assumes that the injection bias is added uniformly across token positions, and
more targeted position-specific injection strategies may shift the ceiling.

\paragraph{Future work.}
Immediate priorities include: (i) sweeping $K > K_\text{max}$ to complete the
ceiling validation; (ii) profiling additional model families (Llama, Mistral,
Gemma) to test $K$-interface transferability; (iii) integrating norm profiling
into training-time monitoring pipelines as in \citet{chen2025persona} to detect
$S_l$ collapse; and (iv) sparse-feature-aware injection \citep{elhage2022toy}
that targets only feature directions most associated with the behavioral circuit.

% ─────────────────────────────────────────────────────────────
\section{Conclusion}
\label{sec:conclusion}
% ─────────────────────────────────────────────────────────────

We have presented \textbf{ksteer}, a calibrated activation steering library that
formalizes injection magnitude using two physically grounded quantities: $K_l$,
the per-dimension RMS of the residual stream, and $S_l$, the behavioral contrast
norm from contrastive prompt pairs.  From these we derive a closed-form coherence
ceiling $K_\text{max}$ and expose a dimensionless integer interface $K$ whose
units are consistent across model families.  On Qwen2.5-3B-Instruct with the
ClearHarm CBRN benchmark, ksteer achieves a 31.5\% relative ASR reduction at
$K_\text{opt} = 4 < K_\text{max} = 6$, with zero gibberish throughout the safe
operating range.  Above $K_\text{max}$, the gibberish rate rises monotonically
from 0 to 0.084 by $K = 9$, confirming the ceiling as a predictive quantity.
Behavioral vectors can be baked into MLP biases for zero-overhead deployment.
We release the ksteer library to facilitate reproducible activation steering
research.

% ─────────────────────────────────────────────────────────────
% Bibliography
% ─────────────────────────────────────────────────────────────
\bibliographystyle{plainnat}
\bibliography{ksteer}

\end{document}
